\section{Reversing maps on bounded tuples}
\label{found:section-reversal}

Throughout this Section, $\K$ is categorical unboundedly often,
$\ell=\LS(\K)$, and $E$ is the presentation of
Theorem~\ref{found:order-template}. The aim is to reverse a strong
self-map on any one fixed bounded tuple in every sufficiently large
categoricity model. We first construct index orders with uniformly
controlled short-tuple orbits.  Homogeneous-order orbit bounds already
play a role in Shelah's work; compare \cite[IV.5.1]{Shelah2009}.
The alignment and elimination below establish the form needed at arbitrary
categoricity cardinals.

For comparison, \cite[Fact~1.4 and the discussion of Claim~1.5,
pp.~3--4]{BoneyVasey2015} records syntactic-characterization results
under exponentiation or cofinality conditions and a gap in the attempted
unrestricted extension. The bounded-tuple reversal needed here is proved
directly in Theorem~\ref{found:reversal}.

\subsection{Finite-support orders and their cuts}

\begin{definition}
\label{found:orders-definition}
For an infinite limit ordinal $\delta$, set
\[
\Gamma_\delta=\Z^{(\delta)},\qquad
H(\delta)=\Q^{(\Gamma_\delta)}.
\]
Parentheses denote finite support. Both groups are ordered by their
greatest differing coordinate. Write $e_\alpha$ for the integer basis
vector at $\alpha<\delta$ and $u_g$ for the rational basis vector at
$g\in\Gamma_\delta$. The rank of a nonzero vector is its greatest
support coordinate. Unless group operations are explicitly mentioned,
$H(\delta)$ is used as a pure linear order; in particular
$\Aut(H(\delta))$ means its order automorphism group.

A cut of an order $L$ is a partition $(L^-,L^+)$ with $L^-<L^+$.
It is an interior nonprincipal cut when both sides are nonempty,
$L^-$ has no greatest element, and $L^+$ has no least element. Its
two characters are $\cf(L^-)$ and $\ci(L^+)$. It is balanced if
these cardinals are equal. Cofinality and coinitiality here refer
to the indicated linear orders.
\end{definition}

\begin{lemma}
\label{found:order-basic}
The order $H(\delta)$ is dense without endpoints and has cardinality
$|\delta|$. Its points have character $\cf(\delta)$ on both sides.
All nonempty bounded open intervals have one order type. All left
rays have one order type, and all right rays have one order type.
An increasing map between ordinal position sets induces an order
embedding between the corresponding $\Gamma$ and $H$ orders by
moving positions and padding by zero. Each point of $H(\delta)$ has
a code using finitely many ordinal positions and integer and
rational coefficients; comparisons are determined by these discrete
coefficients and the order/equality pattern of the ordinal positions.
\end{lemma}

\begin{proof}
Finite-support counting gives
$|\Gamma_\delta|=|H(\delta)|=|\delta|$. The group
$\Gamma_\delta$ has least positive element $e_0$, so every $g$ has
predecessor $g-e_0$. Its cofinality and coinitiality are both
$\cf(\delta)$. Indeed, the vectors $e_\alpha$ for $\alpha$ in a
cofinal subset of $\delta$ are cofinal, and their negatives are
coinitial. The supports of fewer than $\cf(\delta)$ vectors are
bounded in $\delta$, so such a family is neither cofinal nor
coinitial.

The order $H(\delta)$ is a divisible nonzero ordered abelian group,
and hence is dense and has no endpoints. If $(g_i)$ is coinitial
in $\Gamma_\delta$, then $(u_{g_i})$ is coinitial in the positive
cone of $H(\delta)$: given a positive vector of rank $g$, choose
$g_i<g$. Conversely, the ranks of fewer than $\cf(\delta)$ positive
vectors have a strict lower bound in $\Gamma_\delta$, and a positive
basis vector at that lower rank lies below all of them. Thus the
positive cone has coinitiality $\cf(\delta)$. Translation and
negation give that character at every point on both sides.

Translations move points and therefore identify rays of a fixed
direction. We show that any positive vector can be sent to any
other by an increasing automorphism. First translate the index
group $\Gamma_\delta$ so that their greatest support coordinates
agree, say at $g$. Write the two vectors as
\[
a u_g+v,\qquad b u_g+w,
\qquad a,b\in\Q_{>0},\quad \rk(v),\rk(w)<g,
\]
allowing zero lower parts. Fix all other basis vectors and replace
$u_g$ by
\[
\frac ba u_g+\frac{w-v}{a}.
\]
This is an invertible triangular rational-linear map. It preserves
the sign of every nonzero vector, since it either fixes the greatest
nonzero coordinate or multiplies its coefficient by $b/a>0$.
It sends the first displayed vector to the second. Combining this
map with translations sends any ordered pair of distinct points to
any other such pair. Restricting gives the interval assertion.

Finally, an increasing map on ordinal positions preserves the
greatest differing position of two integer vectors. The induced
map on these integer vectors preserves the greatest differing
position of rational vectors as well. This proves the embedding
claim. Expanding the finitely many integer-vector positions used
by a rational vector gives the asserted finite code and comparison
rule.
\end{proof}

\begin{lemma}[Balance of cuts]
\label{found:balanced-cuts}
Every interior nonprincipal cut of $\Gamma_\delta$, and every such
cut of $H(\delta)$, is balanced.
\end{lemma}

\begin{proof}
First consider a cut $(L,R)$ of $\Gamma_\delta$. Among the ranks of
$r-l$, with $l\in L$ and $r\in R$, choose the least ordinal
$\alpha$. Fix one crossing pair of that rank and its common head
above $\alpha$. In that head, each fiber with a fixed integer
coefficient at $\alpha$ lies wholly on one side: a crossing pair
within a fiber would have smaller rank. The resulting cut in the
integer coefficients separates two adjacent integers. A point
with a head above the fixed head would exceed the fixed upper
point, and a point with a head below it would precede the fixed
lower point. The two bordering fibers are therefore cofinal
and coinitial toward the original cut, and their tails are
copies of $\Z^{(\alpha)}$. Their end
characters agree by negation in that tail group. The case
$\alpha=0$ would give a last lower and a first upper point, and
is excluded. This proves balance in $\Gamma_\delta$.

Now let $(L,R)$ be an interior nonprincipal cut of $H=H(\delta)$,
and put
\[
T=\{\rk(r-l):l\in L,\ r\in R\}\subseteq\Gamma_\delta.
\]
This is a final segment: from a crossing pair of rank $g$, subtract
a positive basis vector at any $h>g$ from its lower point. The new
pair still crosses and has rank $h$.

Suppose first that $T$ has a least member $g$. Fix a crossing head
above $g$. Every fiber with a fixed rational coefficient at $g$
lies on just one side of the cut. Otherwise a crossing difference
would have rank below $g$. The induced cut of $\Q$ is nonempty
and proper. If neither coefficient side has a nearest rational,
countable rational approximations are cofinal and coinitial in
the original cut. If the lower coefficients have a greatest
member, the entire fiber at that coefficient lies on the lower
side and is cofinal there. Its tail, supported below $g$, has
greatest available rank $g-e_0$, and the vectors
$n u_{g-e_0}$ and $-n u_{g-e_0}$ show that both its end characters
are $\aleph_0$. On the upper side, rational coefficients decreasing
to the greatest lower coefficient give coinitiality $\aleph_0$.
If instead the upper coefficients have a least member, the
reversed calculation gives cofinality and coinitiality
$\aleph_0$. These cases exhaust the rational cuts, so the original
cut is balanced.

Suppose that $T$ has no least member, and let
$\epsilon=\ci(T)$. Put $S=\Gamma_\delta\setminus T$ and let
$H_S$ be the subgroup supported on $S$. We prove that the lower
character is $\epsilon$ by separating two cases.

Assume there is $l\in L$ for which
\[
\{\rk(r-l):r\in R\}
\quad\text{is coinitial in }T.
\]
The coset $l+H_S$ is contained in $L$, since a crossing within it
would have rank in $S$. It is cofinal in $L$. In fact, if a point
$l'\in L$ lay above the entire coset, $l'-l$ would have positive
rank $h\in T$. Choose $r\in R$ with $\rk(r-l)<h$. Then
$r-l<l'-l$, contradicting $r>l'$. The set $S$ is nonempty, since
otherwise the cofinal coset would be the singleton $\{l\}$.
It has no greatest member: if $s$ were greatest, $s+e_0$ would
be the least member of $T$. Balance in $\Gamma_\delta$, applied
to $(S,T)$, gives $\cf(S)=\epsilon$. The cofinality of $H_S$ at
its upper end equals $\cf(S)$. A cofinal set of ranks gives
cofinal positive basis vectors, while fewer ranks are bounded
and cannot give a cofinal family of vectors. Hence
$\cf(L)=\epsilon$ in this case.

In the other case, each $l\in L$ has a bound $t_l\in T$ strictly
below all ranks $\rk(r-l)$, $r\in R$. Choose crossing pairs
$(l_i,r_i)$, $i<\epsilon$, whose ranks are coinitial in $T$.
For fixed $l\in L$, choose $i$ with
$\rk(r_i-l_i)\leq t_l$. If $l_i\leq l$, then
$0<r_i-l\leq r_i-l_i$, whence
$\rk(r_i-l)\leq\rk(r_i-l_i)\leq t_l$, a contradiction.
Thus $(l_i)_{i<\epsilon}$ is cofinal in $L$.
If fewer than $\epsilon$ points of $L$ were cofinal, their bounds
$t_l$ would have a common strict lower bound in $T$, since
$\epsilon=\ci(T)$. A crossing pair below that bound would have
its lower point above all those points by the same calculation.
This is impossible, proving $\cf(L)=\epsilon$.

Apply the proved lower-side calculation after negating the order
and exchanging $L$ and $R$. Negation preserves every difference
rank, so it has the same $T$ and $\epsilon$. It yields
$\ci(R)=\epsilon$. This proves balance also in the no-least-rank
case.
\end{proof}

\subsection{The orbit record of a short sequence}

For a subset $A$ of a linear order $L$ and a cut
$A=A^-\sqcup A^+$, including cuts with an empty side, define its
\emph{strict slot} by
\[
S_L(A^-,A^+)=\{x\in L\setminus A:A^-<x<A^+\}.
\]
Its occupancy is the assertion that this set is nonempty.
The order/equality diagram of a sequence identifies its distinct
entries and the cuts of their induced order, so the occupancy
data can be indexed without choosing an enumeration of the cuts.

\begin{lemma}[Segments approaching a small cut]
\label{found:cut-segments}
Let $L$ be a dense linear order without endpoints, all of whose
points have character $\kappa$ on both sides. Suppose all bounded
open intervals have one order type and every interior nonprincipal
cut is balanced. For every infinite regular $\epsilon<\kappa$,
the following segment types depend only on $\epsilon$ and the
indicated direction:
\[
[p,\mathfrak c)=\{x\in L^-_{\mathfrak c}:p\leq x\},
\qquad
(\mathfrak c,p]=\{x\in L^+_{\mathfrak c}:x\leq p\},
\]
where $\mathfrak c$ is an interior nonprincipal cut of balanced
character $\epsilon$ and $p$ belongs to the appropriate side.
The isomorphisms can be chosen to match the displayed endpoint.
\end{lemma}

\begin{proof}
We prove both directions simultaneously by induction on the
infinite regular cardinal $\epsilon$. Consider two lower-side
segments and choose strictly increasing sequences
$(p_i)_{i<\epsilon}$ and $(q_i)_{i<\epsilon}$ cofinal in their
respective cut sides, with the indicated endpoints at index zero.
At a stage $i<\epsilon$, fewer than $\epsilon$ preceding points
are not cofinal, so one can choose the next point above them and
above the next point of a fixed cofinal sequence. This gives the
required strictly increasing sequences.

Partition the first segment into blocks as follows:
\[
B_0=\{p_0\},\qquad B_{j+1}=(p_j,p_{j+1}].
\]
For a nonzero limit $i<\epsilon$, set
\[
D_i=\{x\in L:(\exists j<i)\ x\leq p_j\},\quad
U_i=L\setminus D_i,\quad
\mathfrak c_i=(D_i,U_i),\quad B_i=(\mathfrak c_i,p_i].
\]
The cut $\mathfrak c_i$ is interior and nonprincipal. Its lower
side has no greatest point and has cofinality $\cf(i)$: the
strictly increasing $i$-sequence is cofinal there, and any
smaller proposed cofinal family would give a cofinal set of
indices in $i$. Its upper side contains $p_i$. If that side had
a least point $u$, the earlier sequence would give
$\cf(L_{<u})=\cf(i)<\kappa$, contradicting point character.
Balance therefore gives $\mathfrak c_i$ character $\cf(i)$ on
both sides. Since $i<\epsilon$, this is a regular cardinal
strictly below $\epsilon$.

Make the corresponding blocks from the $q_i$. Successor blocks
are isomorphic by bounded-interval homogeneity, adjoining their
greatest points. At a nonzero limit index $i$, use the
\emph{upper-side} induction assertion at $\cf(i)$ to match the
blocks. Match the singleton blocks at zero. These are ordered
partitions of the two segments. Indeed, for any point in the
first segment take the least $i$ such that it is at most $p_i$;
the three cases for $i$ put it in exactly the block just
specified. Gluing the block maps proves the lower-side assertion.

For the upper-side assertion choose strictly decreasing
coinitial sequences. The blocks are $\{p_0\}$ and
$[p_{j+1},p_j)$ at successors. At a nonzero limit $i$, put
\[
V_i=\{x\in L:(\exists j<i)\ p_j\leq x\},\qquad
W_i=L\setminus V_i,\qquad \mathfrak d_i=(W_i,V_i).
\]
The set $V_i$ has no least point and has coinitiality
$\cf(i)$. The set $W_i$ contains $p_i$; if it had a
greatest point, that point's upper character would be
$\cf(i)<\kappa$. Thus $\mathfrak d_i$ is nonprincipal,
and balance gives character $\cf(i)$ on both sides.
The block at $i$ is $[p_i,\mathfrak d_i)$, to which the
\emph{lower-side} induction assertion at $\cf(i)$ applies.
The least index with $p_i\leq x$ verifies that these blocks,
in reverse index order, partition the upper-side segment.
For $\epsilon=\aleph_0$ there are no nonzero limit indices below
$\epsilon$, so this simultaneous induction starts with just the
successor blocks. All displayed endpoint points are matched.
\end{proof}

\begin{theorem}[Short-sequence order orbits]
\label{found:order-orbits}
Let $\delta$ be an infinite limit ordinal and let $\theta$ be an
infinite cardinal with $\theta<\cf(\delta)$. Two sequences of
length at most $\theta$ in $H(\delta)$ lie in the same
$\Aut(H(\delta))$-orbit if and only if they have the same
order/equality diagram and the same occupancies of all strict
slots over their distinct entries.
\end{theorem}

\begin{proof}
Necessity follows from preservation of order. For sufficiency,
the order/equality diagrams give an increasing bijection
$f:A\longrightarrow B$ between the sets of distinct entries,
matching the indexed sequences. We extend it separately on
each occupied slot.

Choose a point $p$ in a slot over $A=A^-\sqcup A^+$. The portion
of the slot at or below $p$ has one of three forms. If
$A^-=\varnothing$, it is $(-\infty,p]$. If $A^-$ has greatest
member $a$, it is $(a,p]$. Otherwise put
\[
D=\{x:(\exists a\in A^-)\ x\leq a\},\qquad U=H(\delta)\setminus D.
\]
The cut $(D,U)$ has no greatest lower point, has $p$ on its
upper side, and has lower character
$\cf(A^-)\leq\theta$. Its upper side has no first point,
since such a point would have character at most $\theta$.
By Lemma~\ref{found:balanced-cuts} this is an interior
nonprincipal cut of balanced character $\cf(A^-)$. The
slot portion is $((D,U),p]$.

In the first case its type is fixed by ray homogeneity; in the
second case by bounded-interval homogeneity; in the third case
by Lemma~\ref{found:cut-segments}. The portion at or above $p$
has the corresponding classification using $A^+$, a least
member when present, and otherwise $\ci(A^+)$. The same
descriptions show that the slot has neither endpoint.

The order map $f$ preserves which boundary case occurs and the
cofinality or coinitiality of the relevant part of $A$. Thus
the corresponding occupied slot over $B$ has matching two
halves. Choose an interior point $q$ there and match the two
halves with their common point $p$ sent to $q$. They glue to a
slot isomorphism. The occupancy assumption ensures that empty
slots also correspond. The entry sets and their occupied slots
partition the whole order. Choice for this set of slots and
gluing with $f$ produce an increasing bijection of
$H(\delta)$, as required.
\end{proof}

\subsection{Bounded names for model orbits}

We now pass from automorphisms of the index order to automorphisms of its
model.  At the cofinalities covered by the orbit theorem, a support record
determines a model orbit, but different records may determine the same
orbit.  Only a uniform upper bound on their number is needed for reversal.

Fix an infinite tuple length $\theta$. A presentation of a
$\theta$-tuple in $E(H(\delta))$ specifies its labels and their
finite index lists. Flatten those lists into an indexed sequence
of at most $\theta$ points of $H(\delta)$. Its \emph{presentation
record} consists of the labels, the order/equality diagram of
this flattened sequence, and all its slot occupancies.

\begin{lemma}[Uniform orbit bound]
\label{found:orbit-bound}
There is a fixed set $\mathcal R_\theta$ of possible presentation
records, depending only on $\ell$ and $\theta$, with
\[
|\mathcal R_\theta|\leq B(\ell,\theta):=2^{2^{\ell+\theta}}.
\]
For every categoricity cardinal
$\kappa\geq\max(\ell,\theta^+)$, the set of automorphism orbits of
$\theta$-tuples in its model has cardinality at most
$B(\ell,\theta)$. This applies also when $\kappa$ is singular.
If $|\delta|=\kappa$ and $\cf(\delta)>\theta$, each realized
record in $E(H(\delta))$ determines exactly one such model orbit.
\end{lemma}

\begin{proof}
The finitely many positions belonging to each tuple entry can
be indexed in $\theta\times\omega$. Choices of arities and
labels number at most $2^{\ell+\theta}$. The flattened
order/equality diagrams number at most $2^\theta$; their sets
of cuts have size at most $2^\theta$, and occupancy subsets
number at most $2^{2^\theta}$. The displayed generous bound
therefore bounds a universal record set.

For a fixed $\delta$ with $\cf(\delta)>\theta$, equal records
give an order automorphism matching the flattened lists by
Theorem~\ref{found:order-orbits}. Applying $E$ gives an
automorphism of the model matching the tuples. Hence a
realized record determines one model orbit. Every tuple has
a labeled presentation, so the number of orbits is bounded
by the number of records.

For any $\kappa\geq\max(\ell,\theta^+)$, the ordinal
$\delta=\kappa+\theta^+$, with ordinal addition here, has
cardinality $\kappa$ and cofinality $\theta^+$. The model
$E(H(\delta))$ has cardinality $\kappa$ and is therefore a
copy of the categoricity model. This proves the bound without
any condition on $\cf(\kappa)$.
\end{proof}

Choose a model $C_\kappa$ at each categoricity size needed in a
given set-sized construction. On its $\theta$-tuple orbits put
$p\longrightarrow r$ when some strong self-map sends a
representative of $p$ to a representative of $r$. This relation
is reflexive and transitive: compose maps after using an
automorphism to align the intermediate representatives.
Choose an injection of the orbit set into the fixed cardinal
$B=B(\ell,\theta)$, once for each $\kappa$. These are its
\emph{orbit names}. For any copy of $C_\kappa$, a tuple's
name is independent of the isomorphism used to $C_\kappa$,
because two such isomorphisms differ by an automorphism.

For a suitable $\delta$ of size $\kappa$, let $F_{\kappa,\delta}$
be the partial function from $\mathcal R_\theta$ to $B$ sending
each realized record to its orbit name. The \emph{profile}
of this presentation consists of its named arrow relation
and $F_{\kappa,\delta}$. The profiles belong to a fixed set
of cardinality at most $2^B$; in particular this bound is
independent of $\kappa$ and $\delta$.

\subsection{Preserving every short record in a small subpresentation}

\begin{lemma}[Closure of ordinal supports]
\label{found:ordinal-support-closure}
Let $\theta^+\leq a$ be infinite cardinals and let $\beta$ be
an ordinal of cofinality greater than $a$. Every subset
$S\subseteq\beta$ of size at most $a$ is contained in a set
$T_0\subseteq\beta$ of size at most $a$ with the following
property: for every $P\subseteq T_0$ of size at most $\theta$
and every finite tuple from $\beta$, there is a tuple from
$T_0$ having the same order/equality diagram over $P$.
Moreover $T_0$ is bounded in $\beta$ and has cofinality
greater than $\theta$ in its induced order.
\end{lemma}

\begin{proof}
Construct increasing sets $S_i\subseteq\beta$ for $i<\theta^+$,
starting with $S$. At a successor stage consider all strict
slots over the entire well-ordered set $S_i$. There are at
most $|S_i|+1\leq a$ of them. In each finite slot add every
point; in each infinite slot add a countably infinite subset.
Together with $S_i$ this realizes every finite order/equality
pattern over $S_i$ which occurs in $\beta$. At limits take
unions. The cardinality stays at most $a$ throughout.

Set $T_0=\bigcup_{i<\theta^+}S_i$. Every set
$P\subseteq T_0$ of size at most $\theta$ is contained in
some $S_i$, by regularity of $\theta^+$. For a given finite
tuple, keep any coordinates already in $S_i$ and match its
other coordinates in their respective slots using the next
stage. This proves the asserted property over $P$.
Since $|T_0|\leq a<\cf(\beta)$, $T_0$ is bounded in
$\beta$. Each $S_i$ is bounded too, and its uppermost slot
is infinite. The next stage therefore contains points above
all of $S_i$. Every at-most-$\theta$ subset of $T_0$ is
bounded in this way, proving $\cf(T_0)>\theta$.
\end{proof}

\begin{lemma}[Appending a long indecomposable segment]
\label{found:append-segment}
Under the hypotheses of Lemma~\ref{found:ordinal-support-closure},
suppose $\beta$ is additively indecomposable. Let $T_0$ be as
there, and let $\delta<\beta$ be additively indecomposable,
with $\otp(T_0)<\delta$ and $\cf(\delta)>\theta$.
Place a consecutive segment $J$ of order type $\delta$
strictly after $T_0$ in $\beta$, and set $T=T_0\cup J$.
Then $\otp(T)=\delta$, and $T$ has the same finite-pattern
matching property over its at-most-$\theta$ subsets.
Consequently $H(T)\subseteq H(\beta)$ computes correctly
the occupancies of every strict slot over any sequence of
at most $\theta$ of its points.
\end{lemma}

\begin{proof}
Choose $\alpha<\beta$ above $T_0$. Indecomposability of
$\beta$ gives $\alpha+\delta<\beta$, so
$J=[\alpha,\alpha+\delta)$ is available. Indecomposability
of $\delta$ gives $\otp(T_0)+\delta=\delta$.

Fix a parameter subset $P\subseteq T$ of size at most
$\theta$ and a finite tuple in $\beta$. Match the
coordinates below $\alpha$ into $T_0$ over $P\cap T_0$
using Lemma~\ref{found:ordinal-support-closure}. Keep
coordinates in $J$ fixed. The parameters in $J$, together
with these finitely many retained coordinates, are bounded
in $J$ because $\cf(\delta)>\theta$. Match all coordinates
at or beyond the end of $J$ by suitably many distinct
points of $J$ above that bound, retaining any repetitions.
The three portions stay in their correct mutual order,
and this gives the required finite tuple over $P$.

Here $\Gamma_T=\Z^{(T)}$ is identified with its image in
$\Gamma_\beta$, and $H(T)=\Q^{(\Gamma_T)}$ with its image
in $H(\beta)$. A sequence of at most $\theta$ points of
$H(T)$ uses at most $\theta$ ordinal positions in its
finite codes. A witness in $H(\beta)$ to an occupied
strict slot uses only finitely many additional ordinal
positions. Match those into $T$ over all the parameter
positions and retain the integer and rational
coefficients. Lemma~\ref{found:order-basic} makes every
comparison with the sequence unchanged. The resulting
point of $H(T)$ witnesses that slot. The reverse
implication follows from inclusion.
\end{proof}

\begin{proposition}[Alignment of model orbits]
\label{found:alignment}
Fix the record set and orbit naming at an infinite length
$\theta$. Let $a<\mu$ be categoricity cardinals with
$a\geq\max(\ell,\theta^+)$. Suppose the following data
have one common profile:
\begin{enumerate}
\item a presentation $E(H(\beta))$ of $C_\mu$, where
$\beta$ is additively indecomposable, $|\beta|=\mu$,
and $\cf(\beta)>a$;
\item presentations $E(H(\delta))$ of $C_a$ for a
cofinal subset of $a^+$ of additively indecomposable
ordinals $\delta$ of size $a$ and cofinality greater
than $\theta$.
\end{enumerate}
Then in any copy $N$ of $C_\mu$, every subset of size
at most $a$ is contained in $A\leq_\K N$ of size $a$
such that the orbit name of \emph{every} $\theta$-tuple
in $A$, computed as a tuple of $C_a$, is the same as
its name computed in $N$.
\end{proposition}

\begin{proof}
Identify $N$ with the presentation in the first item.
Represent the requested set by labeled elements and
collect the ordinal positions in the finite codes of
their indices. There are at most $a$ positions. Apply
Lemma~\ref{found:ordinal-support-closure} to obtain
$T_0$. Its order type is below $a^+$. Choose a
profile occurrence $\delta$ from the second item above
$\otp(T_0)$. Since $a^+\leq\cf(\beta)\leq\beta$ and
$\delta<a^+$, Lemma~\ref{found:append-segment} applies
and gives $T$ of order type $\delta$.

The subpresentation $A=E(H(T))$ is strong in $N$ and
has cardinality $a$. It contains the requested set.
For any $\theta$-tuple in $A$, use a labeled
presentation there. Its record is identical in
$E(H(T))$ and $E(H(\beta))$ by the last assertion of
Lemma~\ref{found:append-segment}. The order isomorphism
$\delta\cong T$ identifies its record in the former
with the same record for $E(H(\delta))$. The two
record-to-name functions belong to the common profile
and hence agree. This proves alignment simultaneously
for all $\theta$-tuples of $A$.
\end{proof}

\subsection{The elimination argument}

Alignment compares every short tuple in a small submodel with its image
in a larger model.  We now use repeated profiles at two categoricity sizes
to rule out an orbit arrow with no reverse.  The larger size will absorb
the cost of processing all tuples of the small model.

\begin{lemma}
\label{found:indecomposable-occurrences}
If $\kappa$ is an infinite cardinal and $\tau\leq\kappa$
is infinite regular, the additively indecomposable
ordinals $\delta<\kappa^+$ with $|\delta|=\kappa$ and
$\cf(\delta)=\tau$ are cofinal in $\kappa^+$.
\end{lemma}

\begin{proof}
Given $\eta<\kappa^+$, choose a limit ordinal
$\rho>\eta$ with $\kappa\leq\rho<\kappa^+$ and
$\cf(\rho)=\tau$, for example by appending an ordinal
segment of type $\tau$ beyond $\eta+\kappa$.
Set $\delta=\omega^\rho$. It is additively
indecomposable. A cofinal sequence of exponents gives
a cofinal sequence in $\omega^\rho$ by continuity of
exponentiation. Fewer than $\cf(\rho)$ ordinals below
$\omega^\rho$ have all their Cantor-normal-form
exponents bounded below $\rho$, and are therefore
bounded in $\omega^\rho$. This proves
$\cf(\delta)=\cf(\rho)=\tau$.
Every ordinal below $\omega^\rho$
has a finite Cantor normal form with exponents below
$\rho$. Thus $|\delta|\leq\max(\aleph_0,|\rho|)=\kappa$.
Since $\delta\geq\rho\geq\kappa$ and $\rho>\eta$,
we have $|\delta|=\kappa$, $\delta<\kappa^+$, and
$\delta>\eta$.
\end{proof}

\begin{theorem}[Eventual reversal at categoricity cardinals]
\label{found:reversal}
For each infinite cardinal $\theta$ there is a cardinal
$b_\K(\theta)$ such that, whenever
$\kappa\geq b_\K(\theta)$ is a categoricity cardinal,
$M\in\K$ has size $\kappa$, $f:M\longrightarrow M$
is strong, and $\bar a$ is a tuple of length at most
$\theta$ in $M$, there is a strong
$g:M\longrightarrow M$ with
\[
g(f(a_i))=a_i
\quad\text{for every coordinate }i.
\]
The bound concerns all sufficiently large categoricity
cardinals, of arbitrary cofinality. No equality of the
form $\kappa^\theta=\kappa$ is assumed.
\end{theorem}

\begin{proof}
It suffices to treat length exactly $\theta$; pad a
nonempty shorter tuple by repetitions, and use any
strong self-map for the empty tuple. In a categoricity
model, the asserted reversal is equivalent to symmetry
of the arrow relation on its $\theta$-tuple orbits.
Indeed a reverse orbit arrow can be precomposed and
postcomposed with automorphisms to match the particular
image tuple and original tuple. The converse follows
by applying reversal to a map witnessing an arrow.

Suppose that symmetry fails at unboundedly many
categoricity cardinals. Let $P$ be an infinite cardinal
bounding the profile palette; for example take $P=2^B$.
Choose a strictly increasing list of failing
categoricity cardinals $(\kappa_i)_{i<P^+}$. At each
step choose it large enough that there is an infinite
regular $\tau_i\leq\kappa_i$ exceeding $\ell$, $\theta$,
and every preceding $\kappa_j$, and that
\[
\kappa_i>P,\qquad
\kappa_i>\kappa_j^\theta\quad(j<i).
\]
Such choices are possible by unbounded failure: first
bound the set of preceding cardinals and their powers,
then choose a regular successor cardinal above the
first bounds and a failing cardinal beyond all of
them. This is a set-sized recursion.

At size $\kappa_i$, fix its orbit naming. By
Lemma~\ref{found:indecomposable-occurrences}, the
ordinals of size $\kappa_i$ and cofinality $\tau_i$
used there are cofinal in $\kappa_i^+$. One profile
occurs cofinally often among them: otherwise the
supremum of at most $P<\kappa_i^+$ bounded fibers
would be below the regular cardinal $\kappa_i^+$.
Since there are $P^+$ sizes and at most $P$ profiles,
choose two sizes $a<\mu$ with the same cofinally
recurring profile. They satisfy
\[
\theta^+\leq a,\qquad a^\theta\leq\mu,
\qquad \tau_\mu>a.
\]
Their common named arrow relation has $p\longrightarrow r$
but no $r\longrightarrow p$. Both names occur in
both sizes. Proposition~\ref{found:alignment} applies
inside every $\mu$-sized model, using any of its
cofinal profile occurrences.

We construct $\theta^+$ rounds of an increasing strong
chain of $\mu$-sized models. At the start of round
$i<\theta^+$, choose an aligned $a$-sized submodel
$A_i$ containing all previous $A_j$. This is possible
because their union has cardinality at most $a$.
They are all strong in the current large model, and
coherence will make $A_j\leq_\K A_i$ for $j<i$.

Enumerate all $\theta$-tuples of $A_i$, using
$a^\theta$ steps. At the step for a tuple $\bar u$,
if its current large-model orbit cannot reach $p$,
leave it as it is. If it can reach $p$, compose such
an arrow with $p\longrightarrow r$ and realize the
composite by a strong self-map of the current model.
Regard this map as an inclusion into an isomorphic
extension: rename its codomain so that the image of
each old element becomes that old element. With the
transported identification of the codomain, the
tuple now has name $r$. At every later large stage
its orbit remains unable to reach $p$. Otherwise
transitivity from $r$ would give $r\longrightarrow p$.
An orbit already unable to reach $p$ has the same
persistence property.

Take strong unions at every limit step, including
between rounds. The entire construction has at most
$\theta^+\cdot a^\theta\leq\mu$ many stages in
cardinality. Every large union therefore has size
$\mu$ and is a copy of $C_\mu$ by categoricity.
Every transition between such copies gives an arrow
of the common relation, so the persistence assertion
also holds at limit stages.

The aligned submodels form a strong increasing chain:
for $j<i$, their underlying sets are included and
both are strong in the large model at the start of
round $i$, so coherence applies. Its union
\[
A=\bigcup_{i<\theta^+}A_i
\]
belongs to $\K$, has cardinality $a$, and is a copy
of $C_a$. Any $\theta$-tuple of $A$ is contained
in some $A_i$, since $\cf(\theta^+)>\theta$. After
round $i$ it cannot reach $p$ in a large model.
At the next aligned small model its small orbit
has the same name, and therefore cannot reach $p$
in the common relation. The inclusion of that small
model into $A$ gives an arrow from this earlier
small orbit to its orbit in $A$. If the latter
could reach $p$, transitivity would contradict the
former exclusion. Thus no tuple in $A$ can have
an orbit reaching $p$.

But $A\cong C_a$, so the orbit named $p$ occurs in
$A$; it reaches itself by the identity map. This
contradiction proves that the failing categoricity
cardinals are bounded. A cardinal above their bound
is $b_\K(\theta)$, and the first paragraph gives
the claimed reversals, including all shorter lengths.
\end{proof}
